% =========================
% Appendix A (Supplementary Materials)
% Save as: appendices/appendixA.tex
% =========================

\chapter{Supplementary Materials}
\label{app:supplement}

% -------- Optional local tweaks for listings --------
% Define JSON and YAML lexers for the 'listings' package (not built-in by default).
\lstdefinelanguage{json}{
  basicstyle=\ttfamily\small,
  numbers=left,
  numberstyle=\tiny,
  stepnumber=1,
  numbersep=6pt,
  showstringspaces=false,
  breaklines=true,
  literate=
   *{0}{{{\color{blue}0}}}{1}
    {1}{{{\color{blue}1}}}{1}
    {2}{{{\color{blue}2}}}{1}
    {3}{{{\color{blue}3}}}{1}
    {4}{{{\color{blue}4}}}{1}
    {5}{{{\color{blue}5}}}{1}
    {6}{{{\color{blue}6}}}{1}
    {7}{{{\color{blue}7}}}{1}
    {8}{{{\color{blue}8}}}{1}
    {9}{{{\color{blue}9}}}{1}
    {:}{{{\color{black}{:}}}}{1}
    {,}{{{\color{black}{,}}}}{1}
    {"}{{{\color{red}"}}}{1}
    {{true}}{{{\color{teal}true}}}{4}
    {{false}}{{{\color{teal}false}}}{5}
    {{null}}{{{\color{teal}null}}}{4}
}
\lstdefinelanguage{yaml}{
  basicstyle=\ttfamily\small,
  numbers=left,
  numberstyle=\tiny,
  stepnumber=1,
  numbersep=6pt,
  showstringspaces=false,
  breaklines=true,
  morecomment=[l]{\#},
  morestring=[b]',
  morestring=[b]"
}

% Keep captions small and italic per your guideline (already set in preamble).

\section{Environment Configuration Files}
This section provides example configuration files and environment templates used for local development and deployment. It helps ensure reproducibility and security by avoiding the exposure of sensitive information.
\label{app:env}
% Purpose: Provide safe, reproducible environment templates without exposing secrets.

\subsection*{A.1.1 \texttt{.env.example} (do not commit real secrets)}
This subsection contains a sample .env file, illustrating the required environment variables for running the system locally. Actual secrets should never be committed or shared.
\begin{lstlisting}[language=yaml, caption={Sample .env.example for local development (no secrets)}]
# Server
PORT=8080
NODE_ENV=development
LOG_LEVEL=info

# WebSocket Gateway
WS_MAX_PAYLOAD=1048576
WS_HEARTBEAT_INTERVAL_MS=18000

# Redis (Pub/Sub + cache)
REDIS_HOST=127.0.0.1
REDIS_PORT=6379

# Database
DB_ENGINE=postgres
DB_HOST=127.0.0.1
DB_PORT=5432
DB_USER=app_user
DB_NAME=messaging

# Auth (JWT)
JWT_ISSUER=example.com
JWT_AUDIENCE=messaging-clients
JWT_PUBLIC_KEY_PATH=./keys/jwt.pub

# Metrics
PROMETHEUS_PORT=9090
\end{lstlisting}

\subsection*{A.1.2 Gateway Configuration (YAML)}
Here, a YAML configuration for the WebSocket gateway is provided, showing key parameters for connection limits, timeouts, and rate limiting.
\begin{lstlisting}[language=yaml, caption={WebSocket gateway configuration (sample)}]
gateway:
  idleTimeoutMs: 30000
  maxConcurrentConns: 80000
  maxFrameSize: 262144
  allowOrigins:
    - https://app.example.com
    - https://staging.example.com
rateLimit:
  msgsPerSecond: 30
  burst: 60
  perIpConnLimit: 5
\end{lstlisting}

% Notes:
% - Keep generic values. Real production configs should be documented privately.
% - Avoid hard-coding secrets; prefer SSM/Secrets Manager (see Ch. 6).

% ---------------------------------------------------

\section{Detailed API Specifications (REST \& WebSocket)}
This section documents the main API endpoints and WebSocket event schemas, supporting integration and client development.
\label{app:apispec}

\subsection*{A.2.1 WebSocket Events (Schemas)}
This part describes the structure of WebSocket events exchanged between clients and the server, including message, presence, and typing events.
% Provide JSON examples for request/response frames to complement Ch. 4.
\begin{lstlisting}[language={}, caption={Client -> Server: msg.send}]
{
  "event": "msg.send",
  "data": {
    "channel_id": "c-82f5",
    "content": "Hello, team!",
    "client_ts": 1737112234
  },
  "trace_id": "8e4b0a7b-1c6a-4c5c-9f0e-1b2c3d4e5f6a"
}
\end{lstlisting}

\begin{lstlisting}[language={}, caption={Server -> Client: msg.deliver}]
{
  "event": "msg.deliver",
  "data": {
    "message_id": "m-9ad12",
    "channel_id": "c-82f5",
    "server_ts": 1737112235,
    "sender_id": "u-42bd",
    "content": "Hello, team!"
  },
  "trace_id": "8e4b0a7b-1c6a-4c5c-9f0e-1b2c3d4e5f6a"
}
\end{lstlisting}

\begin{lstlisting}[language={}, caption={Presence and typing events (bidirectional)}]
{
  "event": "presence.update",
  "data": { "user_id": "u-42bd", "state": "online", "last_seen": 1737112237 }
}
\end{lstlisting}

\subsection*{A.2.2 REST Endpoints (Details)}
Here, you will find sample REST API responses for authentication and message retrieval, useful for backend and frontend integration.
\begin{lstlisting}[language={}, caption={POST /v1/auth/token (response)}]
{
  "access_token": "eyJhbGciOiJSUzI1NiIsInR5cCI6IkpXVCJ9...",
  "expires_in": 3600,
  "token_type": "Bearer",
  "scope": "realtime:connect channels:read messages:write"
}
\end{lstlisting}

\begin{lstlisting}[language={}, caption={GET /v1/messages?channel=:id\&page=:n (response)}]
{
  "page": 3,
  "page_size": 50,
  "items": [
    { "message_id": "m-9ad12", "sender_id": "u-42bd", "content": "Hello", "ts": 1737112201 },
    { "message_id": "m-9ad13", "sender_id": "u-51fa", "content": "Hi!",   "ts": 1737112203 }
  ],
  "next_page": 4
}
\end{lstlisting}

\subsection*{A.2.3 Error Model and Rate Limits}
This subsection explains the error response format and rate limiting strategies applied to API endpoints and WebSocket events.
\begin{lstlisting}[language={}, caption={Error envelope (all endpoints/events)}]
{
  "error": {
    "code": "RATE_LIMIT_EXCEEDED",
    "message": "Too many requests",
    "request_id": "r-1f9d22",
    "hint": "Reduce frequency or back off and retry later."
  }
}
\end{lstlisting}

% ---------------------------------------------------


\section{Deployment Scripts (Excerpts)}
Here, you will find excerpts from deployment scripts, including Dockerfiles, docker-compose setups, and Terraform snippets for cloud infrastructure.
\label{app:deploy}

\subsection*{A.7.1 Dockerfile (Gateway)}
This Dockerfile shows how the Node.js gateway service is built and configured for production deployment.
\begin{lstlisting}[language=yaml, caption={Dockerfile (Node.js gateway, simplified)}]
FROM node:20-alpine
WORKDIR /app
COPY package*.json ./
RUN npm ci --omit=dev
COPY . .
ENV NODE_ENV=production
EXPOSE 8080
CMD ["node", "dist/main.js"]
\end{lstlisting}

\subsection*{A.7.2 docker-compose.yml (Local Dev)}
This docker-compose file provides a local development environment with all required services, including Redis and PostgreSQL.
\begin{lstlisting}[language=yaml, caption={docker-compose (local development)}]
version: "3.9"
services:
  redis:
    image: redis:7-alpine
    ports: ["6379:6379"]
  db:
    image: postgres:16-alpine
    environment:
      POSTGRES_DB: messaging
      POSTGRES_USER: app_user
      POSTGRES_PASSWORD: app_pass
    ports: ["5432:5432"]
  gateway:
    build: .
    environment:
      - REDIS_HOST=redis
      - DB_HOST=db
    ports: ["8080:8080"]
    depends_on: [redis, db]
\end{lstlisting}

\subsection*{A.7.3 Terraform Snippet (ALB Listener for WSS)}
This Terraform snippet demonstrates how to configure an AWS Application Load Balancer listener for secure WebSocket traffic.
\begin{lstlisting}[language=yaml, caption={Terraform (HCL) snippet for ALB listener (pseudo)}]
resource "aws_lb_listener" "wss" {
  load_balancer_arn = aws_lb.main.arn
  port              = "443"
  protocol          = "HTTPS"
  ssl_policy        = "ELBSecurityPolicy-TLS-1-2-2017-01"
  certificate_arn   = var.acm_cert_arn
  default_action {
    type             = "forward"
    target_group_arn = aws_lb_target_group.gateway.arn
  }
}
\end{lstlisting}

% ---------------------------------------------------

\section{Experiment Reproducibility Checklist}
This section lists the software versions and step-by-step instructions needed to reproduce the experiments and benchmarks described in the thesis.
\label{app:repro}

\subsection*{A.8.1 Software Versions}
This table details the exact software versions used in all experiments to ensure reproducibility and consistency.
\begin{table}[h]
\centering
\caption{Versions used in experiments}
\label{tab:versions}
\begin{tabular}{l l}
\toprule
\textbf{Component} & \textbf{Version} \\
\midrule
Node.js & 20.x LTS \\
Redis & 7.x \\
PostgreSQL & 16.x \\
k6 & 0.50+ \\
Artillery & 2.x \\
Nginx & 1.25+ \\
\bottomrule
\end{tabular}
\end{table}

\subsection*{A.8.2 Steps to Reproduce (Summary)}
This summary outlines the main steps required to set up the infrastructure, deploy the system, and run performance tests.
\begin{enumerate}
  \item Provision infrastructure (VPC, ALB, compute, ElastiCache, DB).
  \item Deploy gateway container(s) and set environment variables (no secrets in images).
  \item Seed database and create test channels/users.
  \item Run k6/Artillery scenarios from a separate load injector VPC/subnet.
  \item Collect metrics (CloudWatch, Prometheus, logs) and export CSV for analysis.
\end{enumerate}

% ---------------------------------------------------

\section{Additional Logs and Error Samples}
This section provides sample logs and error codes to help with troubleshooting and understanding system behavior during failures.
\label{app:logs}

\subsection*{A.9.1 Structured Log (Masked)}
This log example shows the format and key fields of structured logs generated by the gateway service.
\begin{lstlisting}[language={}, caption={Gateway structured log (JSON)}]
{
  "ts":"2026-01-17T14:20:01.120Z",
  "level":"info",
  "msg":"delivered",
  "channel":"c-82f5",
  "message_id":"m-9ad12",
  "latency_ms":42,
  "conn_id":"k3r2",
  "trace_id":"8e4b0a7b-1c6a-4c5c-9f0e-1b2c3d4e5f6a"
}
\end{lstlisting}

\subsection*{A.9.2 WebSocket Close Codes (Excerpt)}
This table lists common WebSocket close codes and their meanings, useful for diagnosing connection issues.
\begin{table}[h]
\centering
\caption{Common WebSocket close codes (excerpt)}
\label{tab:ws-codes}
\begin{tabular}{l l p{8cm}}
\toprule
\textbf{Code} & \textbf{Name} & \textbf{Meaning} \\
\midrule
1000 & Normal Closure & Regular closure without error \\
1006 & Abnormal Closure & No close frame; network break or peer lost \\
1008 & Policy Violation & Server closed due to policy (e.g., ACL) \\
1011 & Internal Error & Unexpected server condition \\
\bottomrule
\end{tabular}
\end{table}

% ---------------------------------------------------

\section{Glossary of Technical Terms}
This glossary defines important technical terms used throughout the appendix and thesis, supporting reader understanding.
\label{app:glossary}
% Keep concise, alphabetically if possible.
\begin{itemize}
  \item \textbf{Backpressure:} Mechanisms to prevent producers from overwhelming consumers by applying flow control.
  \item \textbf{Fan-out:} Broadcasting a message to multiple recipients (e.g., all members of a channel).
  \item \textbf{Presence:} Representation of a user's online/offline/away state and last-seen timestamp.
  \item \textbf{Sharding:} Partitioning data or load across multiple nodes to scale horizontally.
  \item \textbf{Sticky Session:} Routing a client to the same backend node across requests or connections.
  \item \textbf{Trace ID:} Unique identifier propagated across services for distributed tracing.
\end{itemize}

% ================= End Appendix A =================